\documentclass[11pt]{article}
\usepackage[T1]{fontenc}
\usepackage[utf8]{inputenc}
\usepackage[a4paper,margin=28mm]{geometry}
\usepackage{amsmath,amssymb,booktabs,threeparttable,hyperref}
\hypersetup{colorlinks=true,linkcolor=black,urlcolor=blue,citecolor=black}
\newcommand{\keywords}[1]{\par\noindent\textbf{Keywords:} #1\par}
\title{Productivity and investment under uncertainty}
\author{Jordan Lee\\\small Department of Economics}
\date{Working paper --- \today}
\begin{document}
\maketitle
\begin{abstract}
This economics working-paper starter demonstrates a compact title page,
abstract, JEL-style keywords, equations, a regression table, footnotes, and
references. Replace the illustrative content and follow the target journal's
submission instructions.
\end{abstract}
\keywords{investment; productivity; uncertainty}
\noindent\textbf{JEL codes:} D24; E22; O47.
\section{Introduction}
Explain the economic question, its policy or business relevance, and the
contribution relative to the literature. Keep the identification strategy
visible to readers who are not specialists in the application.
\section{Empirical framework}
Consider the linear specification
\begin{equation}
 y_{it} = \alpha + \beta x_{it} + \gamma_i + \delta_t + \varepsilon_{it},
 \label{eq:regression}
\end{equation}
where fixed effects and the error structure should be defined for the actual
study. Explain the source of variation and the estimand before reporting it.
\section{Results}
\begin{table}[ht]
\centering
\caption{Illustrative regression results.}
\begin{threeparttable}
\begin{tabular}{lcc}
\toprule
 & (1) Baseline & (2) Fixed effects \\
\midrule
Input $x$ & 0.42 & 0.37 \\
 & (0.08) & (0.10) \\
Observations & 240 & 240 \\
\bottomrule
\end{tabular}
\begin{tablenotes}\small
\item Standard errors are illustrative. Add clustering, controls, and fixed
effects notes required by the application.
\end{tablenotes}
\end{threeparttable}
\end{table}
Interpret the coefficient in Equation~\eqref{eq:regression}, discuss robustness,
and separate descriptive results from causal claims.
\section{Conclusion}
State the main result, limitations, and implications for future research.
\begin{thebibliography}{9}
\bibitem{angrist} J. D. Angrist and J.-S. Pischke, \emph{Mostly Harmless
Econometrics}, Princeton University Press, 2009.
\bibitem{imbens} G. W. Imbens and D. B. Rubin, \emph{Causal Inference},
Cambridge University Press, 2015.
\end{thebibliography}
\end{document}
